\chapter{Tumbuhan adalah Produsen}\label{ch:g6:plants-are-producers}

Timbanglah sebuah pot yang disiram, tanamlah sebutir \omterm{def:g2:from-seed-to-plant:seed}{biji} ek di
dalamnya, lalu kembalilah bertahun-tahun kemudian untuk menimbang
pohon ek mudanya: berkilo-kilogram kayu, kulit kayu, \omterm{def:g1:plants-around-us:parts}{daun} dan \omterm{def:g1:plants-around-us:parts}{akar} ---
sementara tanah di dalam potnya nyaris tidak kehilangan segram pun
(\cref{ex:g4:what-plants-need:tree} pertama kali melontarkan kejutan
ini). Berkilo-kilogram pohon, yang dibuat dari apa? Bab ini
memastikan jawabannya --- dan menyebutkan perdagangan yang khas milik
\omterm{def:g1:plants-around-us:plant}{tumbuhan}.

\section{Materi organik dan materi mineral}

\begin{definition}[Materi mineral dan materi organik]\label{def:g6:plants-are-producers:matter}
Materi, bagi seorang ahli \omterm{def:g1:living-or-not:living}{biologi}, ada dua macam:
\begin{itemize}
  \item \emph{materi mineral}\index{materi mineral}: materi dunia yang
        tak hidup --- air, gas di udara, batuan, \omterm{prop:g4:what-plants-need:needs}{garam mineral}
        terlarut pada
        \cref{prop:g4:what-plants-need:needs};
  \item \emph{materi organik}\index{materi organik}: materi yang
        dibuat \omterm{def:g6:exploring-our-environment:components}{makhluk hidup} dan menyusun tubuhnya --- kayu, \omterm{def:g1:plants-around-us:parts}{daun},
        \omterm{def:g3:bones-and-muscles:muscle}{otot}, lemak, gula, wol. Dapat dibakar, dapat melapuk, dan
        selalu dapat dilacak sampai ke seorang pembuat yang hidup.
\end{itemize}
\end{definition}

\begin{example}[Mengelompokkan sesaku benda]\label{ex:g6:plants-are-producers:sorting}
Dari sebuah saku mantel sesudah jalan-jalan: sebutir kerikil ---
mineral; sebutir \omterm{prop:g3:flowers-fruits-seeds:transform}{buah} hazel --- organik (sebuah pohon hazel yang
membuatnya); sebuah kancing plastik --- anehnya organik menurut
asalnya (plastik dibuat dari minyak bumi, yaitu sisa \omterm{def:g6:exploring-our-environment:components}{makhluk hidup}
purba yang sudah berubah); setetes air hujan --- mineral; sebuah sarung
tangan wol --- organik (seekor domba yang membuat wolnya).
\end{example}

\section{Perdagangan para produsen}

\begin{proposition}[Tumbuhan menghasilkan materi organik dari materi mineral]\label{prop:g6:plants-are-producers:produce}
Sebuah \omterm{def:g1:plants-around-us:plant}{tumbuhan} hijau, kalau hanya diberi bahan \emph{mineral} ---
air dan garam terlarut dari tanah, sebuah gas yang diambil dari
\emph{udara} --- serta \omterm{prop:g3:balanced-diet:jobs}{energi} \emph{\omterm{def:g6:exploring-our-environment:conditions}{cahaya}}, membuat materi
\emph{organik}: gulanya sendiri, dan dari gula itu kayu, \omterm{def:g1:plants-around-us:parts}{daun} serta
\omterm{def:g1:plants-around-us:parts}{akar}. Pembuatan inilah perdagangan yang dijanjikan
\cref{prop:g4:what-plants-need:make} untuk dijelaskan: \omterm{def:g1:plants-around-us:plant}{tumbuhan} adalah
\emph{produsen}\index{produsen} --- satu-satunya pembuat \omterm{def:g6:plants-are-producers:matter}{materi
organik} baru dari pasokan yang murni mineral. Zat hijau \omterm{def:g1:plants-around-us:parts}{daunnya} adalah
permesinan penangkap \omterm{def:g6:exploring-our-environment:conditions}{cahayanya}; pembuatannya terjadi di dalam \omterm{def:g1:plants-around-us:parts}{daun},
pada siang hari.
\end{proposition}
\admitted

\begin{omfigure}
\begin{tikzpicture}[scale=0.92]
  % leaf in the middle
  \draw[thick, omMeth, fill=omMeth!20] (5.2,1.5)
    .. controls (5.6,2.4) and (6.8,2.5) .. (7.2,1.6)
    .. controls (6.8,0.8) and (5.6,0.75) .. (5.2,1.5);
  \draw[omMeth!60!black] (5.35,1.5) -- (7.05,1.58);
  \node[font=\small\bfseries, omMeth!60!black] at (6.2,2.75)
    {daun hijaunya, pada siang hari};
  % inputs
  \node[font=\small, align=center] (water) at (1.4,0.35)
    {air dan garam mineral\\ (dari tanah, lewat akarnya)};
  \node[font=\small, align=center] (gas) at (1.4,1.7)
    {sebuah gas dari udara};
  \node[font=\small, align=center] (light) at (1.4,3.1)
    {energi cahaya};
  \draw[->, thick, omDef] (water.east) -- (5.3,1.2);
  \draw[->, thick, omDef] (gas.east) -- (5.15,1.5);
  \draw[->, thick, omProp] (light.east) -- (5.5,2.05);
  % outputs
  \node[font=\small, align=center] (sugar) at (10.6,1.9)
    {\textbf{materi organik}:\\ gula, lalu kayu,\\ daun, akar, biji};
  \draw[->, very thick, omMeth!60!black] (7.3,1.7) -- (sugar.west);
  % oxygen small note
  \node[font=\footnotesize, align=center] (oxy) at (10.6,0.5)
    {(dan oksigen, yang dilepaskan\\ ke udara --- lihat catatannya)};
  \draw[->, omExo] (7.25,1.2) -- (oxy.west);
\end{tikzpicture}

{\small Perdagangan para \omterm{def:g5:food-webs:roles}{produsen}: masukan mineral dan \omterm{def:g6:exploring-our-environment:conditions}{cahaya} masuk,
\omterm{def:g6:plants-are-producers:matter}{materi organik} keluar. Hanya \omterm{def:g1:plants-around-us:plant}{tumbuhan} hijau (dan kerabat hijaunya)
yang menjalankan pembuatan ini.}
\end{omfigure}

\begin{example}[Buktinya, yang dirangkai]\label{ex:g6:plants-are-producers:evidence}
Tiga pengamatan sebelumnya sekarang saling mengunci. Pot yang
ditimbang: berkilo-kilogram pohon eknya tidak datang dari tanahnya ---
kebanyakan datang dari udara dan air. Si tawanan pucat pada
\cref{ex:g4:what-plants-need:pale}: tanpa \omterm{def:g6:exploring-our-environment:conditions}{cahaya}, tanpa pembuatan ---
\omterm{def:g1:plants-around-us:plant}{tumbuhannya} kelaparan dalam gelap dengan gudang mineral yang penuh.
\omterm{def:g1:plants-around-us:parts}{Daun} mint yang kelaparan mineral pada
\cref{exo:g4:what-plants-need:11}: garam memang dibutuhkan --- tetapi
sebagai pasokan yang kecil, bukan sebagai bagian terbesarnya. Bagian
terbesarnya dari udara dan air, \omterm{prop:g3:balanced-diet:jobs}{energinya} dari \omterm{def:g6:exploring-our-environment:conditions}{cahaya}, bumbunya dari
tanah.
\end{example}

\begin{remark}[Pertukaran gas, ke dua arah]\label{rem:g6:plants-are-producers:gas}
Gas udara yang dimasukkan \omterm{def:g1:plants-around-us:parts}{daunnya} adalah \emph{karbon dioksida} ---
gas yang persis sama yang diembuskan setiap makhluk yang bernapas
(\cref{prop:g4:breathing:exchange}); dan pembuatannya melepaskan
\emph{oksigen} ke udara --- gas yang persis sama yang dibutuhkan
setiap makhluk yang bernapas. \omterm{def:g1:plants-around-us:plant}{Tumbuhan} pun bernapas, siang dan malam,
seperti semua \omterm{def:g6:exploring-our-environment:components}{makhluk hidup}; tetapi pada siang hari pembuatan oleh
\omterm{def:g1:plants-around-us:parts}{daun} yang bekerja jauh melampaui \omterm{def:g4:breathing:movements}{pernapasannya}. Udara yang dapat
dihirup di dunia ini adalah hasil sampingan para \omterm{def:g5:food-webs:roles}{produsen} --- satu
utang lagi yang ditanggung setiap \omterm{def:g1:animals-around-us:animal}{hewan} kepada yang hijau.
\end{remark}

\section{Produsen memberi makan semua orang --- dan menyimpan untuk dirinya}

\begin{proposition}[Mata rantai pertama, yang dijelaskan]\label{prop:g6:plants-are-producers:firstlink}
\Cref{prop:g3:food-chains:start} menetapkan bahwa setiap \omterm{def:g3:food-chains:chain}{rantai
makanan} bermula pada sebuah \omterm{def:g1:plants-around-us:plant}{tumbuhan}; sekarang ketetapan itu punya
alasannya. Hanya \omterm{def:g5:food-webs:roles}{produsen} yang membuat \omterm{def:g6:plants-are-producers:matter}{materi organik} dari \omterm{def:g6:plants-are-producers:matter}{materi
mineral}; setiap \omterm{def:g5:food-webs:roles}{konsumen} (\cref{def:g5:food-webs:roles}) hanya dapat
\emph{mengubah} \omterm{def:g6:plants-are-producers:matter}{materi organik} yang dimakannya. Jadi seluruh \omterm{def:g6:plants-are-producers:matter}{materi
organik} dalam setiap tubuh yang hidup --- kelinci, rubah, elang, kamu
--- pertama kali dibuat di dalam sehelai \omterm{def:g1:plants-around-us:parts}{daun}, dan piramida makanan
pada \cref{prop:g5:food-webs:pyramid} berdiri di atas lantai hijaunya
karena di situlah materi memasuki dunia yang hidup.
\end{proposition}
\admitted

\begin{example}[Gudang milik tumbuhannya sendiri]\label{ex:g6:plants-are-producers:pantry}
Apa yang dibuat sebuah \omterm{def:g1:plants-around-us:plant}{tumbuhan} melebihi kebutuhan harinya, disimpannya
sebagai cadangan organik: \omterm{ex:g4:plants-without-seeds:bulbtuber}{umbi batang} kentang dan \omterm{ex:g4:plants-without-seeds:bulbtuber}{umbi lapis} tulip
(\cref{ex:g4:plants-without-seeds:bulbtuber}), kedua belahan tebal
kacang (\cref{def:g2:from-seed-to-plant:seed}), \omterm{def:g1:plants-around-us:parts}{akar} wortel yang gemuk
(\cref{exo:g1:plants-around-us:10}), rasa manis \omterm{prop:g3:flowers-fruits-seeds:transform}{buah} bit.
Masing-masingnya adalah \omterm{def:g6:exploring-our-environment:conditions}{cahaya} musim panas, yang ditabung sebagai
materi --- dan masing-masingnya persis yang dipanen petani, dan yang
disantap pemakannya.
\end{example}

\begin{method}[Melacak makanan apa pun sampai ke daunnya]\label{met:g6:plants-are-producers:trace}
Ambil makanan apa pun di atas meja:
\begin{enumerate}
  \item kalau ia materi \omterm{def:g1:plants-around-us:plant}{tumbuhan} --- gandum pada rotinya, kentangnya,
        apelnya --- ia langsung merupakan buatan \omterm{def:g5:food-webs:roles}{produsen}: sebutkan
        \omterm{def:g1:plants-around-us:plant}{tumbuhannya} dan organ simpanannya;
  \item kalau ia materi \omterm{def:g1:animals-around-us:animal}{hewan} --- keju, \omterm{def:g2:animals-grow-up:born}{telur}, ham --- turunlah
        menyusuri rantainya: \omterm{def:g1:animals-around-us:animal}{hewannya} makan apa? lalu siapa pembuat
        makanan itu? teruskan sampai ke \omterm{def:g1:plants-around-us:plant}{tumbuhannya};
  \item tuliskan rantainya dengan anak panah
        (\cref{def:g3:food-chains:chain});
  \item kesimpulannya selalu sama: jejaknya berakhir di sehelai \omterm{def:g1:plants-around-us:parts}{daun},
        di dalam \omterm{def:g6:exploring-our-environment:conditions}{cahaya}.
\end{enumerate}
\end{method}

\section{Latihan}

\begin{exercise}[$\star$]\label{exo:g6:plants-are-producers:1}
Definisikan \omterm{def:g6:plants-are-producers:matter}{materi mineral} dan \omterm{def:g6:plants-are-producers:matter}{materi organik}, dengan dua contoh
masing-masing.
\end{exercise}

\begin{exercise}[$\star$]\label{exo:g6:plants-are-producers:2}
Nyatakan perdagangan para \omterm{def:g5:food-webs:roles}{produsen}: masukannya, sumber \omterm{prop:g3:balanced-diet:jobs}{energinya},
keluarannya, dan bengkel tempat ia terjadi.
\end{exercise}

\begin{exercise}[$\star$]\label{exo:g6:plants-are-producers:3}
Kelompokkan sebagai mineral atau organik: embun, \omterm{def:g1:animals-around-us:covering}{cangkang} siput, lilin
lebah, granit, minyak goreng, garam laut.
\end{exercise}

\begin{exercise}[$\star$]\label{exo:g6:plants-are-producers:4}
Dari mana kebanyakan kilogram pohon eknya datang, kalau nyaris bukan
dari tanahnya?
\end{exercise}

\begin{exercise}[$\star$]\label{exo:g6:plants-are-producers:5}
Gas mana yang diambil \omterm{def:g1:plants-around-us:parts}{daun} yang bekerja dari udara, dan gas mana yang
dilepaskannya? Siapa yang mengembuskan gas yang pertama dan
membutuhkan gas yang kedua?
\end{exercise}

\begin{exercise}[$\star$]\label{exo:g6:plants-are-producers:6}
Sebutkan empat organ simpanan \omterm{def:g1:plants-around-us:plant}{tumbuhan} dan apa yang disimpan
masing-masing, dari \cref{ex:g6:plants-are-producers:pantry}.
\end{exercise}

\begin{exercise}[$\star\star$]\label{exo:g6:plants-are-producers:7}
Mengapa \omterm{def:g1:plants-around-us:plant}{tumbuhan} di dalam lemari yang gelap kelaparan padahal gudang
\omterm{prop:g4:what-plants-need:needs}{garam mineralnya} penuh? Masukan mana yang tidak punya pengganti?
\end{exercise}

\begin{exercise}[$\star\star$]\label{exo:g6:plants-are-producers:8}
Mengapa tidak ada \omterm{def:g1:animals-around-us:animal}{hewan} yang dapat menjadi mata rantai pertama sebuah
\omterm{def:g3:food-chains:chain}{rantai makanan}? Jawablah dengan pembedaan \omterm{def:g5:food-webs:roles}{produsen}/\omterm{def:g5:food-webs:roles}{konsumen}.
\end{exercise}

\begin{exercise}[$\star\star$]\label{exo:g6:plants-are-producers:9}
Jalankan \cref{met:g6:plants-are-producers:trace} atas sebuah roti isi
keju: rotinya, menteganya, kejunya --- tiga jejak menuju tiga \omterm{def:g1:plants-around-us:parts}{daun}.
\end{exercise}

\begin{exercise}[$\star\star$]\label{exo:g6:plants-are-producers:10}
``Kompos memberi makan \omterm{def:g1:plants-around-us:plant}{tumbuhannya} sebagaimana jerami memberi makan
seekor sapi.'' Betulkan kalimat itu dengan tepat, memakai
\cref{rem:g4:what-plants-need:soil} dan definisi dalam bab ini --- apa
yang sesungguhnya disediakan kompos, dan apa yang dibuat \omterm{def:g1:plants-around-us:plant}{tumbuhannya}
untuk dirinya sendiri?
\end{exercise}

\begin{exercise}[$\star\star$]\label{exo:g6:plants-are-producers:11}
Sebuah akuarium tertutup yang tersinari, berisi \omterm{def:g1:plants-around-us:plant}{tumbuhan} air dan \omterm{prop:g3:sorting-living-things:groups}{ikan},
dapat berjalan berbulan-bulan; sebuah akuarium tertutup yang
\emph{gelap} gagal dalam hitungan hari. Pakai
\cref{rem:g6:plants-are-producers:gas} untuk menjelaskan kedua
hasilnya.
\end{exercise}

\begin{exercise}[$\star\star\star$]\label{exo:g6:plants-are-producers:12}
Sebuah percobaan tua: sebatang anakan pohon dedalu yang ditumbuhkan
lima tahun di dalam sebuah bak bertambah berat berkilo-kilogram; tanah
kering di dalam baknya hanya kehilangan beberapa gram. Sang pelaku
percobaan menyimpulkan ``pohonnya terbuat dari air belaka''. Bagian
kesimpulannya yang mana yang benar, yang mana yang belum lengkap ---
dan masukan mineral kedua apa, yang tidak diketahuinya, yang akan
ditambahkan seorang ahli \omterm{def:g1:living-or-not:living}{biologi} modern?
\end{exercise}

\section{Soal: Setahun Perjalanan Petani Sayur}

\begin{problem}[{Soal akhir pekan --- cahaya, daun dan pembukuan
sebuah kebun sayur}]\label{pb:g6:plants-are-producers:1}
Seorang petani sayur menanam kentang, wortel, kacang dan selada, dan
menyimpan pembukuan yang tajam tentang apa yang masuk dan apa yang
keluar.

\medskip\noindent\textbf{Bagian I --- Teka-teki pembukuannya.}
\begin{enumerate}
  \item Kebunnya menghasilkan dua ton sayur setahun; petaninya hanya
        menebarkan beberapa karung kompos. Timbanglah teka-tekinya:
        dari mana datangnya ton-ton itu, kalau bukan dari karungnya?
        (Jawablah dengan perdagangan para \omterm{def:g5:food-webs:roles}{produsen}.)
  \item Kelompokkan hasil panennya: kentangnya, wortelnya, \omterm{def:g2:from-seed-to-plant:seed}{biji}
        kacangnya, \omterm{def:g1:plants-around-us:parts}{daun} seladanya --- sebutkan peran masing-masing
        bagi \omterm{def:g1:plants-around-us:plant}{tumbuhannya} (organ simpanan, organ yang bekerja, simpanan
        \omterm{def:g2:from-seed-to-plant:seed}{biji}).
  \item Petaninya berkata: ``pemasokku yang sesungguhnya mengantar
        pada siang hari.'' Terjemahkan ke dalam istilah
        \cref{prop:g6:plants-are-producers:produce}.
  \item Mengapa petaninya tetap membutuhkan komposnya, kalau bukan
        untuk ton-tonnya? (Ingatlah \omterm{def:g1:plants-around-us:parts}{daun} mint di dalam pasir.)
\end{enumerate}

\medskip\noindent\textbf{Bagian II --- Membaca musimnya.}
\begin{enumerate}[resume]
  \item Selada bulan Juni, yang dinaungi kain percobaan selama
        seminggu, keluar dalam keadaan pucat dan memanjang. Sebutkan
        percobaan yang diulanginya, dan masukan yang disingkirkannya.
  \item Barisan kentangnya ditimbuni tanah dan bagian atasnya yang
        rimbun dipuji: ``bagian atas yang baik, umbi yang baik.''
        Jelaskan pepatahnya --- apa hubungan antara \omterm{def:g1:plants-around-us:parts}{daun} yang bekerja
        di atas dan simpanan yang membengkak di bawah?
  \item Kemarau bulan Agustus menghentikan pertumbuhannya walaupun
        mataharinya menyala-nyala. Masukan mineral mana yang gagal,
        dan mengapa \omterm{def:g6:exploring-our-environment:conditions}{cahaya} tidak dapat menggantikannya?
  \item Pada bulan September petaninya memotong lalu menyingkirkan
        tanaman kacangnya tetapi membenamkan \omterm{def:g1:plants-around-us:parts}{akarnya} ke dalam tanah.
        Seorang tetangga menyebut ini pemborosan. Belalah kebiasaannya
        dengan definisi materinya --- apa yang sedang dikembalikan ke
        dalam tanah, untuk siapa?
\end{enumerate}

\medskip\noindent\textbf{Bagian III --- Di luar kebunnya.}
\begin{enumerate}[resume]
  \item Babi milik petaninya memakan kentang yang kecil-kecil;
        keluarganya memakan babinya. Tuliskan rantainya, lalu nyatakan
        di mana \omterm{def:g6:plants-are-producers:matter}{materi organiknya} semula dibuat.
  \item Seorang pembeli mengaku bahwa saladnya ``kebanyakan matahari''.
        Auditlah pengakuannya: daftarkan bahan yang sesungguhnya pada
        \omterm{def:g1:plants-around-us:parts}{daun} saladnya dan peran matahari yang tepat.
  \item Embun beku mematikan selada yang terakhir dalam semalam;
        keduanya melapuk di tempatnya menjelang bulan November. \omterm{def:g6:plants-are-producers:matter}{Materi
        organiknya} tidak lenyap. Dahuluilah
        \cref{ch:g6:decomposers-and-soil}: siapa yang mengambilnya,
        dan menuju keadaan apa?
  \item Rangkumlah kebunnya dalam satu kalimat yang memuat: \omterm{def:g5:food-webs:roles}{produsen},
        \omterm{def:g6:plants-are-producers:matter}{materi mineral}, \omterm{def:g6:plants-are-producers:matter}{materi organik}, \omterm{def:g6:exploring-our-environment:conditions}{cahaya}.
\end{enumerate}
\end{problem}
